\section{Addressing and collision statistics}\label{app:addressing}

Two practical questions arise once stabilized windows are read as outputs of an overlap transducer: how long does it take to forget the initial state, and how likely is an address collision among independently sampled objects?

\subsection{Overlap transducer and exact synchronization}\label{sec:causality:sync}

\begin{definition}[Overlap transducer]\label{def:transducer}
Fix $m\ge 2$. Define the deterministic finite-state transducer
\[
\widehat{\mathcal{T}}_m=(V_m,\{0,1\},\hat{\delta}_m,\hat{\mathrm{out}}_m),
\]
with state set $V_m:=\{0,1\}^{m-1}$. For a state $u=(u_1,\dots,u_{m-1})\in V_m$ and an input bit $c\in\{0,1\}$, set
\[
\hat{\delta}_m(u,c):=(u_2,\dots,u_{m-1},c),
\qquad
\hat{\mathrm{out}}_m(u,c):=\Fold_m(u_1,\dots,u_{m-1},c)\in X_m.
\]
\end{definition}

\begin{definition}[Synchronization time]\label{def:sync-time}
For an input sequence $a=(a_1,a_2,\dots)\in\{0,1\}^{\NN}$, define
\[
T_m(a):=\inf\Bigl\{n\ge 1:\hat{\delta}_m^n(u,a_{1:n})=\hat{\delta}_m^n(v,a_{1:n})\ \text{for all }u,v\in V_m\Bigr\},
\]
with the convention $\inf\emptyset=\infty$.
\end{definition}

\begin{theorem}[Exact synchronization]\label{thm:exact-sync}
For every $m\ge 2$ and every input sequence $a\in\{0,1\}^{\NN}$ one has
\[
T_m(a)=m-1.
\]
\end{theorem}

\begin{proof}
Fix $u=(u_1,\dots,u_{m-1})\in V_m$ and $n<m-1$. Iterating the shift-register transition gives
\[
\hat{\delta}_m^n(u,a_{1:n})=(u_{n+1},\dots,u_{m-1},a_1,\dots,a_n).
\]
Hence the state after $n$ steps still contains the initial component $u_{n+1}$, so choosing two initial states $u,v$ that differ at index $n+1$ yields $\hat{\delta}_m^n(u,a_{1:n})\neq \hat{\delta}_m^n(v,a_{1:n})$. Therefore $T_m(a)\ge m-1$.

For $n=m-1$, the same identity gives $\hat{\delta}_m^{m-1}(u,a_{1:m-1})=(a_1,\dots,a_{m-1})$, which is independent of $u$. Thus the synchronization condition holds at $n=m-1$ and $T_m(a)\le m-1$.
\end{proof}

\subsection{Collision-limited addressing after deterministic burn-in}\label{sec:causality:capacity}

\begin{definition}[Input law, address distribution, and collision probability]\label{def:collision}
Let $\nu_u$ be a shift-invariant probability measure on input sequences $\{0,1\}^{\NN}$, and let $\nu_u^{(m)}$ denote its length-$m$ marginal. The induced address distribution on stabilized types is
\[
\nu_{u,m}:=(\Fold_m)_*\nu_u^{(m)}\in\Delta(X_m),
\qquad
\nu_{u,m}(x)=\nu_u^{(m)}(\mathcal{F}_m(x)).
\]
Its collision probability is
\[
\Col_u(m):=\sum_{x\in X_m}\nu_{u,m}(x)^2=\exp(-h_2(\nu_{u,m})),
\]
where $h_2(\nu):=-\log\sum_x \nu(x)^2$ is the R\'enyi entropy of order $2$.
\end{definition}

\begin{remark}
In the geometric source setting $\nu_u$ is the symbolic law induced by a cut-and-project scan. The theorem below is formulated for general shift-invariant $\nu_u$ to separate the protocol-level content from the source law.
\end{remark}

\begin{theorem}[Exact occupancy identity after deterministic burn-in]\label{thm:exact-occupancy}
Fix $m\ge 2$ and let $\nu_u$ be a shift-invariant input law on $\{0,1\}^{\NN}$. Consider $N$ independent objects, each producing an independent input sequence $A^{(i)}\sim\nu_u$. Fix a burn-in budget $n\ge m-1$ and define the post-burn-in $m$-address of object $i$ by
\[
X^{(i)}_{m,n}:=\Fold_m(A^{(i)}_{n-m+2:n+1})\in X_m.
\]
Let $\mathsf{Fail}_{m,n,N}$ denote the event that there exists an address collision among $\{X^{(i)}_{m,n}\}_{i=1}^N$. Then
\begin{equation}\label{eq:exact-occupancy}
\PP\bigl(X^{(1)}_{m,n},\dots,X^{(N)}_{m,n}\text{ all distinct}\bigr)=N!\,e_N(\nu_{u,m}),
\end{equation}
where
\[
e_N(\nu_{u,m})
:=
\sum_{\substack{S\subset X_m\\|S|=N}}
\prod_{x\in S}\nu_{u,m}(x)
\]
is the $N$-th elementary symmetric polynomial of the address weights. Equivalently,
\begin{equation}\label{eq:exact-failure}
\PP(\mathsf{Fail}_{m,n,N})=1-N!\,e_N(\nu_{u,m}).
\end{equation}
\end{theorem}

\begin{proof}
By Theorem~\ref{thm:exact-sync}, every object is synchronized once $n\ge m-1$, so the post-burn-in address of each object has law $\nu_{u,m}$. Because the objects have independent input sequences and each address is a measurable function of one object's input, the random variables
\[
X^{(1)}_{m,n},\dots,X^{(N)}_{m,n}
\]
are i.i.d.\ with common law $\nu_{u,m}$.

Therefore
\[
\PP\bigl(X^{(1)}_{m,n}=x_1,\dots,X^{(N)}_{m,n}=x_N\bigr)=\prod_{i=1}^N \nu_{u,m}(x_i).
\]
Hence
\[
\PP\bigl(X^{(1)}_{m,n},\dots,X^{(N)}_{m,n}\text{ all distinct}\bigr)
=
\sum_{\substack{(x_1,\dots,x_N)\in X_m^N\\x_i\neq x_j\ \forall i\neq j}}
\prod_{i=1}^N \nu_{u,m}(x_i).
\]
Group this sum by the unordered $N$-element subsets
\[
S=\{x_1,\dots,x_N\}\subset X_m.
\]
For each fixed $S$ with $|S|=N$, there are exactly $N!$ orderings of its elements, and every ordering contributes the same product $\prod_{x\in S}\nu_{u,m}(x)$. Therefore
\[
\PP\bigl(X^{(1)}_{m,n},\dots,X^{(N)}_{m,n}\text{ all distinct}\bigr)
=
N!\sum_{\substack{S\subset X_m\\|S|=N}}
\prod_{x\in S}\nu_{u,m}(x)
=
N!\,e_N(\nu_{u,m}),
\]
which proves \eqref{eq:exact-occupancy}. Equation \eqref{eq:exact-failure} is the complementary probability.
\end{proof}

\begin{corollary}[Collision-limited addressing after burn-in]\label{cor:collision-budget}
In the setting of Theorem~\ref{thm:exact-occupancy},
\begin{equation}\label{eq:collision-budget}
\PP(\mathsf{Fail}_{m,n,N})\le \binom{N}{2}\Col_u(m).
\end{equation}
\end{corollary}

\begin{proof}
The union bound gives
\[
\PP(\mathsf{Fail}_{m,n,N})\le \sum_{1\le i<j\le N}\PP(X^{(i)}_{m,n}=X^{(j)}_{m,n}).
\]
Each pair is i.i.d.\ with common law $\nu_{u,m}$, so
\[
\PP(X^{(i)}_{m,n}=X^{(j)}_{m,n})=\sum_{x\in X_m}\nu_{u,m}(x)^2=\Col_u(m).
\]
Summing over pairs proves \eqref{eq:collision-budget}.
\end{proof}

For the uniform raw-word law, all occupancy quantities can be written directly in terms of the fiber counts. We record the corresponding formulas in the next subsection as a benchmark case for the general theory.
